
\chapter{12.5 用坐标表示平移}


\subsection*{一、选择题（共 8 小题）}

\questionlist

\item 在平面直角坐标系中，若将三角形上各点的纵坐标都减去 3，横坐标保持不变，则所得图形在原图形基础上（\quad ）

  \optionlist
  \item 向左平移了 3 个单位
  \item 向下平移了 3 个单位
  \item 向上平移了 3 个单位
  \item 向右平移了 3 个单位
  \end{enumerate}

\item 在平面直角坐标系中，$\triangle ABC$ 的三个顶点坐标分别为 $A(4, 5)$，$B(1, 2)$，$C(4, 2)$，将 $\triangle ABC$ 向左平移 5 个单位后，$A$ 的对应点 $A_1$ 的坐标是（\quad）

  \optionlist
  \item (0, 5)
  \item (-1, 5)
  \item (9, 5)
  \item (-1, 0)
  \end{enumerate}

\item 点 $A(4, 3)$ 经过平移后得点 $B(6, -3)$，它的平移过程是（\quad）

  \optionlist
  \item 向右平移 2 个单位后再向下平移 6 个单位
  \item 向左平移 2 个单位后再向下平移 2 个单位
  \item 向左平移 2 个单位后再向上平移 6 个单位
  \item 向右平移 6 个单位后再向上平移 2 个单位
  \end{enumerate}

\item 在平面直角坐标系中，已知点 $O(0, 0)$，$A(2, 4)$，将线段 $OA$ 沿 $x$ 轴向左平移 2 个单位，再向上平移 2 个单位，得到 $O, A$ 两点的对应点分别为 $O_1, A_1$，则 $O_1, A_1$ 的坐标分别是（\quad）

  \optionlist
  \item (-2, 2), (4, 6)
  \item (-2, 2), (0, 6)
  \item (2, -2), (4, 6)
  \item (2, -2), (0, 6)
  \end{enumerate}

\item 在平面直角坐标系中，将点 $A(x, y)$ 向左平移 5 个单位长度，再向上平移 3 个单位长度后与点 $B(-3, 2)$ 重合，则点 $A$ 的坐标是（\quad）

  \optionlist
  \item (2, 5)
  \item (-8, 5)
  \item (-8, -1)
  \item (2, -1)
  \end{enumerate}

\item 在 $\triangle ABC$ 内任意一点 $P(a, b)$ 经过平移后对应点 $P_1(c, d)$，已知 $A(3, 2)$ 在经过此次平移后对应点 $A_1$ 的坐标为 $(5, -1)$，则 $a+b-c-d$ 的值为（\quad）

  \optionlist
  \item -5
  \item -1
  \item 1
  \item 5
  \end{enumerate}

\item 在坐标平面上两点 $A(-a+2, -b+1)$，$B(3a, b)$，若点 $A$ 向右移动 2 个单位长度后，再向下移动 3 个单位长度后与点 $B$ 重合，则点 $B$ 所在的象限为（\quad）

  \optionlist
  \item 第一象限
  \item 第二象限
  \item 第三象限
  \item 第四象限
  \end{enumerate}






\item 图，已知直角坐标系中的点 $A, B$ 的坐标分别为 $A(2, 4)$，$B(4, 0)$，且 $P$ 为 $AB$ 的中点。若将线段 $AB$ 向右平移 3 个单位后，与点 $P$ 对应的点为 $Q$，则点 $Q$ 的坐标是（\quad）
  \par\smallskip\noindent
  \begin{minipage}[t]{0.44\linewidth}
  \vspace{0pt}
 \optionlist
  \item (3, 2)
  \item (6, 2)
  \item (6, 4)
  \item (3, 5)
  \end{enumerate}
  \end{minipage}
  \hfill
  \begin{minipage}[t]{0.2\linewidth}
  \questionimage{12.5-8.png}
  \end{minipage}
  


\end{enumerate}

\subsection*{二、填空题（共 6 小题）}

\blanklist[9]

\item 如图，将线段 $AB$ 平移，使 $B$ 点到 $C$ 点，则平移后 $A$ 点的坐标为\_\_\_\_\_\_。

% 插入图片：此处应有第9题图
\rightquestionimage[width=0.35\textwidth]{0.35\textwidth}{12.5-9.png}


\item 在平面直角坐标系中，有一条线段 $AB$，已知点 $A(-2, 0)$ 和 $B(0, 2)$，平移线段 $AB$ 得到线段 $A_1B_1$。若点 $A$ 的对应点 $A_1$ 的坐标为 $(1, 3)$，则线段 $A_1B_1$ 的中点坐标是\_\_\_\_\_\_。
\vspace{1cm}

\item 已知点 $P(-4, 6-3b)$，先向左平移 2 个单位，再向下平移 3 个单位，恰好落在 $x$ 轴的负半轴上，则 $b$ 应为\_\_\_\_\_\_。
\vspace{2cm}
\item 在平面直角坐标系中，规定一个点先向上平移 2 个单位，再向右平移 1 个单位为 1 次运动。点 $(-2, -3)$ 经过\_\_\_\_\_\_次这样的运动后到达点 $P'(4, 9)$。
\vspace{1cm}
\item 在平面直角坐标系中，孔明玩走棋的游戏，其走法是：棋子从原点出发，第 1 步向右走 1 个单位长度，第 2 步向右走 2 个单位长度，第 3 步向上走 1 个单位长度，第 4 步向右走 1 个单位长度……依此类推，第 $n$ 步的走法是：当 $n$ 能被 3 整除时，则向上走 1 个单位长度；当 $n$ 能被 3 整除时，则向下走 1 个单位长度；当 $n$ 能被 3 整除时，若 $n$ 能被 3 整除，则向右走 2 个单位长度，若 $n$ 能被 3 整除，则向左走 2 个单位长度。当走完第 100 步时，棋子所在位置的坐标是\_\_\_\_\_\_。
\vspace{2cm}
\item 初一年级某班有 54 名学生，所在教室有 6 行 9 列座位，用 $(m, n)$ 表示第 $m$ 行第 $n$ 列的位置，新学期准备调整座位，设某个学生原来的座位为 $(m, n)$，如果调整后的座位为 $(i, j)$，则称该生作了平移 $(a, b) = [m-i, n-j]$，并称 $a+b$ 为该生的位置数。若某生的位置数为 10，则当 $m+n$ 取最小值时，$m+n$ 的最大值为\_\_\_\_\_\_。
\vspace{2cm}

\end{enumerate}

\subsection*{三、解答题（共 4 小题）}

\solutionlist[15]

\item 五边形 $ABCDE$ 的顶点坐标分别为 $A(0, 6)$、$B(-3, -3)$、$C(-1, 0)$、$D(1, 0)$、$E(3, 3)$，将五边形 $ABCDE$ 看成经过一次平移后得到 $A_1B_1C_1D_1E_1$，其中顶点 $A$ 的对应点是 $A_1(-3, 10)$。

(1) 请写出其它对应点的坐标；

(2) 请写出这一平移的平移方向和平移距离。

\vspace{4cm}

\item 已知点 $P(2a-12, 1-a)$ 位于第三象限，点 $Q(x, y)$ 位于第二象限且是由点 $P$ 向上平移一定单位长度得到的。

(1) 若点 $P$ 的纵坐标为 $-3$，试求出 $a$ 的值；

(2) 在 (1) 的条件下，试求出符合条件的点 $Q$ 的坐标。

若点 $P$ 的横、纵坐标都是整数，试求出 $a$ 的值以及线段 $PQ$ 长度的取值范围。

\rightquestionimage[width=0.35\textwidth]{0.35\textwidth}{12.5-16.png}

\vspace{5cm}


\item 图中的图形是将坐标 $(0, 0), (1, 0), (3, 0), (2, 1), (3, 4), (5, 3), (5, 2), (3, 2)$ 的点用线段依次连接而成的。

% 插入图片：此处应有第17题图

下面将以上各点做如下变化：

(1) 横坐标保持不变，纵坐标分别乘以 $-1$，所得图案与原图案有什么变化？

(2) 横坐标和纵坐标都乘以 $-1$，所得图案与原图案相比有什么变化？

(3) 横坐标加 $1$，纵坐标加 $2$，所得图案与原图案相比有什么变化？

\rightquestionimage[width=0.35\textwidth]{0.35\textwidth}{12.5-17.png}

\vspace{5cm}

\newpage 

\item 操作与探究：

(1) 对数轴上的点 $P$ 进行如下操作：先把点 $P$ 表示的数乘以 $\frac{1}{3}$，再把所得数对应的点向右平移 $1$ 个单位，得到点 $P'$ 的对应点 $P'$。

点 $A, B$ 在数轴上，对线段 $AB$ 上的每个点进行上述操作后得到线段 $A'B'$，其中点 $A, B$ 的对应点分别为 $A', B'$。如图1，若点 $A$ 表示的数是 $-3$，则点 $A'$ 表示的数是\_\_\_\_\_\_；若点 $B'$ 表示的数是 $2$，则点 $B$ 表示的数是\_\_\_\_\_\_；已知线段 $AB$ 上的点 $E$ 经过上述操作后得到的对应点 $E'$ 与点 $E$ 重合，则点 $E$ 表示的数是\_\_\_\_\_\_。

% 插入图片：此处应有图1

(2) 如图2，在平面直角坐标系 $xOy$ 中，对正方形 $ABCD$ 及其内部的每个点进行如下操作：把每个点的横、纵坐标都乘以同一个实数 $a$，将得到的点先向右平移 $m$ 个单位，再向上平移 $n$ 个单位（$m > 0, n > 0$），得到正方形 $A'B'C'D'$ 及其内部的点，其中点 $A, B$ 的对应点分别为 $A', B'$。已知正方形 $ABCD$ 内部的一个点 $F$ 经过上述操作后得到的对应点 $F'$ 与点 $F$ 重合，求点 $F$ 的坐标。

% 插入图片：此处应有图2
\rightquestionimage[width=0.35\textwidth]{0.35\textwidth}{12.5-18.png}

\vspace{5cm}

\end{enumerate}
